\section{面对新模型与新缩写，怎么判断它值不值得学？}
\label{intro:23e-tech-judgement}

我们已从语言模型走到系统工具、部署，再回看 RoPE、MoE、GQA 等架构。
技术名称不断变化，但读者不必为了追赶每个缩写而重新背一遍名词表。
更重要的是拿到一个新技术时，能独立判断\textbf{它解决了哪个瓶颈，
用什么证据证明，付出了什么代价，以及哪些条件下会失效}。

\subsection{先把“改进”分到五层，而不是混在一起}

\begin{enumerate}
\item \textbf{表示层：}Tokenizer、Embedding、位置表示如何编码输入？
改动会影响形状、长度或位置感知，不能直接等同于回答更准确。
\item \textbf{模型计算层：}Attention、FFN、MoE、归一化与缓存改变什么？
区分训练 FLOPs、解码延迟、KV 内存和总参数量。
\item \textbf{训练与反馈层：}预训练、SFT、偏好优化、RLVR 是否修改权重？
数据、奖励与评测目标一致吗？
\item \textbf{推理策略层：}采样、验证、搜索、上下文示例是否使用额外计算，
但通常保持同一权重？
\item \textbf{系统层：}RAG、工具、MCP、Workflow、Agent/A2A 是怎样连接
资料与授权动作的？安全责任归谁？
\end{enumerate}
例如“模型接入了新的报表搜索工具”，首先是\textbf{系统与证据来源}变化；
不能说这证明基础模型的 Transformer 宽度更大。

\subsection{Scaling Laws：更多参数一定更聪明吗？}

经典缩放规律研究在一定实验范围内发现：
语言模型的\textbf{验证损失}常随参数量、训练 Token 与计算预算
呈较规律的变化趋势。Kaplan 等（2020）讨论了这些关系，
见 \href{https://arxiv.org/abs/2001.08361}{Scaling Laws for Neural Language Models}。
Hoffmann 等（2022）进一步研究给定训练计算量时
\textbf{模型参数和训练 Token 应如何分配}，见
\href{https://arxiv.org/abs/2203.15556}{Training Compute-Optimal Large Language Models}。

这里的词语是\textbf{经验规律}，不是“参数翻倍就必然答对某道题”的定理。
参数、数据、算力与质量互相制约；
评测分布、推理预算以及是否有工具，也会改变最终结果。
规模有价值，但不能替代任务定义、数据质量和可靠性验证。

\subsection{Bitter Lesson：计算与通用方法为何常被重视？}

Richard Sutton 在
\href{http://www.incompleteideas.net/IncIdeas/BitterLesson.html}{The Bitter Lesson}
中主张，长期来看，能够借助增长的计算规模进行学习或搜索的通用方法，
往往比大量手写领域规则更有扩展性。
这是一种关于研究历史与方法选择的\textbf{经验观点}，
不是针对每一个数据量、预算和应用场景的严格证明。
在受限硬件、特定安全规则和明确工作流中，
手写约束与专家知识仍然可以非常重要。
第一部分问“为什么规则难以写尽”，在这里得到了一次回响：
\textbf{让系统学会一般模式，与让应用遵守明确规则，并不矛盾。}

\subsection{拿到宣传语，先做一次拆解}

假设听到一句：“新注意力让上下文扩大十倍，速度提升十倍，
因此所有业务的答案都会更准确。”
先要求补齐缺失的证据：
\begin{itemize}
\item \textbf{比较对象：}与哪一个基线、哪一种实现、哪一个硬件平台比较？
\item \textbf{指标定义：}更快指首 Token 延迟、吞吐还是总运行时间？
上下文变长不自动等于能准确找到远处证据。
\item \textbf{质量代价：}长程检索、普通对话、代码等任务是否都测过？
有没有关键回归？
\item \textbf{实验条件：}输入长度、并发数、精度、缓存策略和训练量是否相同？
\item \textbf{应用边界：}引用核查、工具权限、失败处理有没有另外验证？
\end{itemize}
一项新技术可能主要减少显存，
另一项主要增加训练期活跃参数，
第三项主要减少人为操作；\textbf{三者不能只按同一个“更先进”排名}。

\subsection{把全书压缩成五道可以反复使用的问题}

读完任何新的架构、训练算法或 Agent 系统，试着回答：
\begin{enumerate}
\item 上一个方法\textbf{哪一步不够用}，能否构造反例？
\item 新方法的\textbf{输入、输出、形状或调用协议}是什么？
\item 有什么可复算的\textbf{数字、测试或可追溯来源}？
\item 哪一项\textbf{成本或风险}被转移到别处？
\item 什么情形下应该\textbf{退回更简单的方法}？
\end{enumerate}

\par\medskip\noindent{\sffamily\footnotesize\color{blue}\textbf{读完自测}}\quad
\phantomsection\label{check:23e-tech-judgement}
如果一种新方法把 KV Cache 元素减少一半，是否就证明回答质量一定提升？
至少还要测哪两类指标？
\quad{\footnotesize\hyperref[answer:23e-tech-judgement]{核对解答。}}\par\medskip

\begin{conclusion}
从标量训练到 Transformer，再到工具系统，核心方法始终是
\textbf{明确问题、给出可运行的运算或接口、检查结果与代价}。
面对新缩写，不妨先判断它作用在哪一层，再决定是否需要学它的细节。
\end{conclusion}
